\SourceLeaf{452}%
\begin{center}
\begin{tabular}{@{}l p{.54\linewidth} r@{}}
\(a.\)&\raggedright Länge von 1505 Messstangen ohne Korrektur \(.\ .\ .\ .\ .\ .\ .\ .\ .\)&\(3010{,}00000\) Toisen\\
\(b.\)&\raggedright Summe der mit den Glaskeilen gemessenen Abstände und der zu berücksichtigenden ganzen und halben Durchmesser der Ablotungszylinder \(.\)&\(+3{,}58389\) T.\\
\(c.\)&\raggedright Länge der Ergänzungsstange \(.\ .\ .\ .\ .\ .\ .\ .\ .\ .\ .\ .\ .\ .\ .\)&\(+1{,}22106\) T.\\
\(d.\)&\raggedright Korrektur für die Neigung der Ablotungszylinder gegenüber der Lotlinie~\(.\)&\(-0{,}00008\) T.\\
\(e.\)&\raggedright Korrektur für die Abweichung der Stangen von der Fluchtlinie \(.\ .\ .\ .\ .\)&\(-0{,}00051\) T.\\
\(f.\)&\raggedright Korrektur für die fehlerhaften Längen der Messstangen \(.\ .\ .\ .\ .\ .\ .\)&\(-0{,}10245\) T.\\
\(g.\)&\raggedright Korrektur für die Abweichung der Temperatur der Messstangen von \(13^{0}\) R.&\(-0{,}19906\) T.\\
\(h.\)&\raggedright Reduktion auf die Meeresoberfläche \(.\ .\ .\ .\ .\ .\ .\ .\ .\ .\ .\)&\(-0{,}02264\) T.\\\cline{3-3}
&\raggedright Länge der Braacker Basis nach der neuen Berechnung \(.\ .\ .\ .\ .\)&\(=3014{,}48021\) Toisen
\end{tabular}
\end{center}

Doch auch diese Berechnung hat einen Schwachpunkt, nämlich die unter \(g\) aufgeführte Korrektur für die Temperatur der Messstangen. Eine Beschreibung des bei der Basismessung verwendeten Apparats mit Abbildungen hat \textsc{Schumacher} in der Schrift ‘Schreiben an Dr. \textsc{Olbers} in Bremen etc. etc.’, Altona 1821’, veröffentlicht. Daraus geht hervor, dass die Temperaturen nicht mit Metallthermometern, sondern mit gewöhnlichen, eingelegten Thermometern bestimmt wurden. Das ist schon für sich genommen ein misslicher Umstand. Noch viel schwerer wiegt aber, dass die Ausdehnbarkeit der Stangen nur aus einigen abends und morgens im Feld vorgenommenen Messungen ihrer Längen abgeleitet wurde. Dem lässt sich jedoch abhelfen. Im Jahr 1853 wurde nämlich die Stange Nr. IV. des Basisapparats von \textsc{Schumacher} nach Pulkowa gebracht, um sie direkt mit den dort gesammelten Maßnormalen zu vergleichen. Bei dieser Gelegenheit wurde auch die Ausdehnung dieser Stange für \(100^{0}\) ermittelt. Berechnet man den von \textsc{Struve} (siehe ‘Arc du méridien entre le Danube et la mer glaciale’, S. 51) angegebenen Wert der Ausdehnungskoeffizienten, so erhält man für die Korrektur unter \(g\) \(-0{,}22812\) statt \(-0{,}19906\).

Mit dieser Berichtigung, die auch Professor \textsc{Peters} übernimmt, erhält man dann für die Länge der \textit{Braacker Basis}
\[
=3014{,}451\ \text{Toisen},
\]
und dieser Wert ist als \textit{endgültig} anzusehen. Ich füge hinzu, dass die Angabe in diesen Toisen auf dem Vergleich mit dem Pulkowaer \textsc{Fortin}-Normal beruht. Da dieses aber bis auf eine verschwindende Kleinigkeit mit der \textsc{Bessel}schen Toise übereinstimmt, kann die Länge auch ohne Weiteres als in \textsc{Bessel}schen Toisen ausgedrückt gelten.’

Das obige Koordinatenverzeichnis ergibt für die Länge der Basis \(5875{,}3614\) der dort verwendeten Einheiten oder \(3014{,}5757\) Toisen bei einem Erdmeridian von \(360\times57009{,}746\) Toisen.

\textsc{Eduard Schmidt. Gauss} an \textsc{Schumacher}: Göttingen 1830 April 30. ‘Um Ihr Vertrauen in \textsc{Schmidt}s Berechnung zu stärken, bemerke ich, dass er die beiden Hauptelemente der Erddimensionen viermal berechnet hat. — — Das Ergebnis (IV) hat er mir handschriftlich mitgeteilt. Es ist dasselbe, das meinen neuen Hilfstafeln zugrunde liegt, nämlich die Abplattung \(\frac{1}{297{,}732}\) und \(\frac{\text{Erdquadrant}}{90}=57008^{\mathrm{T}}551\).’

\textsc{Bessel}. ‘Ueber einen Fehler in der Berechnung der französischen Gradmessung und seinen Einfluss auf die Bestimmung der Figur der Erde.’ Astronomische Nachrichten Nr. 438, Band 19. Seite 116. 1841 Dezember 2. ‘Mittlerer Grad des Meridians \(=57013{,}109\) Toisen, große Halbachse \(a=3272077{,}14\) Toisen, kleine Halbachse \(b=3261139{,}33\) Toisen, \(a:b=299{,}1528:298{,}1528\)’

Bei der Anwendung der in den obigen Verzeichnissen angegebenen Koordinaten sind diese also vorläufig, bis die Basis und die Verbindungsdreiecke bis \textit{Hamburg — Hohenhorn} neu gemessen worden sind, mit dem folgenden Korrekturfaktor zu multiplizieren:
